\subsection{A multiplier criterion for harmonic defining functions}
\label{sec:harmonic-multiplier-criterion}

The following criterion permits multipliers that degenerate at the
singular set. Its spectral application appears in
\cref{sec:spectral-positivity}.

Let \(F\not\equiv0\) be a homogeneous polynomial of positive degree on
\(\R^N\), \(N\ge2\), satisfying \eqref{eq:harmonic-minimal-weight}, and set
\(U=\{S>0\}\).  Let \(h\) be smooth on \(U\), put \(a=e^h\),
and suppose that a smooth function \(\tau\) on \(U\) satisfies
\begin{equation}\label{eq:harmonic-log-compatibility}
 \Gamma(F,h)=\tau F.
\end{equation}
Thus \(\nabla h\) is tangent to the regular zero set. Introduce
\begin{equation}\label{eq:harmonic-log-auxiliary}
 H_h=\Gamma(h,h),\qquad Q_h=H_h+2\tau,\qquad
 P_h=Q_h\Delta h-\tfrac12\Gamma(h,Q_h),
\end{equation}
and the boundary coefficient
\begin{equation}\label{eq:harmonic-log-D0}
 D_h=S(\Delta h-\tau-H_h)-q-\tfrac12\Gamma(h,S).
\end{equation}
These functions are defined on \(U\) without division by \(F\).

\begin{lemma}[Harmonic multiplier identity]
\label{lem:harmonic-log-multiplier}
With the preceding notation, \(G=aF\) satisfies
\begin{align}
 |\nabla G|^2&=a^2(S+F^2Q_h),
 \label{eq:harmonic-log-gradient}\\
 \cL(G)&=a^3F\left(D_h+F^2P_h-
                    \tfrac F2\Gamma(F,Q_h)\right).
 \label{eq:harmonic-log-defect}
\end{align}
Along \(C_{\reg}=\{F=0,S>0\}\), one has
\begin{equation}\label{eq:harmonic-log-Jacobi}
 \frac{D_h}{S}=-\frac{J_Cw}{w},\qquad
 w=(a\sqrt S)^{-1}.
\end{equation}
If in addition \(\Gamma(F,Q_h)=\kappa F\) for a smooth function
\(\kappa\), then \eqref{eq:harmonic-log-defect} becomes
\begin{equation}\label{eq:harmonic-log-two-coefficients}
 \cL(G)=a^3F(D_h+F^2D_{1,h}),\qquad
 D_{1,h}=Q_h\Delta h-\tfrac12\kappa-
                         \tfrac12\Gamma(h,Q_h).
\end{equation}
\end{lemma}

\begin{proof}
The product and chain rules give
\[
 \nabla G=a(\nabla F+F\nabla h),\qquad
 \Delta G=aF(\Delta h+Q_h).
\]
The first formula gives \eqref{eq:harmonic-log-gradient}.  Set
\(V=S+F^2Q_h\).  Differentiating \(|\nabla G|^2=a^2V\) yields
\[
 a^{-3}\Gamma(G,|\nabla G|^2)
 =2FV(\tau+H_h)+\Gamma(F,V)+F\Gamma(h,V).
\]
Here
\[
 \begin{aligned}
 \Gamma(F,V)&=2qF+2FSQ_h+F^2\Gamma(F,Q_h),\\
 \Gamma(h,V)&=\Gamma(h,S)+2\tau F^2Q_h+
                                      F^2\Gamma(h,Q_h).
 \end{aligned}
\]
Substitute these identities in
\(\cL(G)=|\nabla G|^2\Delta G-
\tfrac12\Gamma(G,|\nabla G|^2)\), and use
\(Q_h=H_h+2\tau\).  The terms proportional to \(F\) and
\(F^3\) give \eqref{eq:harmonic-log-defect}.
The additional divisibility assumption gives
\eqref{eq:harmonic-log-two-coefficients}.

On the regular zero set, \(\cK_F(a)=a^3D_h\), so
\eqref{eq:boundary-potential-algebraic} gives
\eqref{eq:harmonic-log-Jacobi}.
\end{proof}

\begin{theorem}[A monotonicity criterion for harmonic multipliers]
\label{thm:harmonic-monotonicity-criterion}
In the setting of \cref{lem:harmonic-log-multiplier}, suppose on
\(U\) that
\begin{equation}\label{eq:harmonic-log-monotonicity}
 Q_h\ge0,\qquad D_h\ge0,\qquad \Delta h\ge0,\qquad
 F\Gamma(F,Q_h)\le0,\qquad \Gamma(h,Q_h)\le0.
\end{equation}
Then \(\nabla(aF)\ne0\) on \(U\), and
\(F\cL(aF)\ge0\) there. With
\(Z=\R^N\setminus U=\{\nabla F=0\}\), the field
\[
 X=\frac{\nabla(aF)}{|\nabla(aF)|}\quad\text{on }U
\]
is a two-sided subcalibration of \(E=\{F<0\}\) across \(Z\).
Consequently \(E\) is locally perimeter minimizing and
\(\partial\llbracket E\rrbracket\) is locally mass minimizing.
No regularity of the spherical link is assumed.
\end{theorem}

\begin{proof}
By \eqref{eq:harmonic-log-gradient},
\(|\nabla(aF)|^2\ge a^2S>0\).  In the bracket in
\eqref{eq:harmonic-log-defect}, each of
\(D_h\), \(F^2Q_h\Delta h\),
\(-F^2\Gamma(h,Q_h)/2\), and
\(-F\Gamma(F,Q_h)/2\) is nonnegative.  Thus
\(F\cL(aF)\ge0\), and
\eqref{eq:div-normalized-gradient} gives the required divergence
signs.  At every regular zero, \(\nabla(aF)=a\nabla F\), so
\(X\) is the exterior unit normal to \(E\).

By \cref{rem:polynomial-finite-perimeter,lem:square-free-critical-set},
\(E\) has locally finite perimeter and its algebraic critical set
\(Z\subset\{F=0\}\) has locally zero \(\cH^{N-1}\)-measure.
Apply \cref{thm:removable-subcalibration} to \(X\); no extension
of \(a\) across \(Z\) is needed.
\end{proof}

\begin{remark}[Alternative coefficient tests]
\label{rem:harmonic-alternative-tests}
If \(\Gamma(F,Q_h)=\kappa F\), the monotonicity conditions may be replaced by
\(Q_h\ge0\), \(D_h\ge0\), and \(D_{1,h}\ge0\).
Without this divisibility assumption, the conditions
\[
 Q_h\ge0,\quad D_h\ge0,\quad P_h\ge0,\quad
 \Gamma(F,Q_h)^2\le16D_hP_h
\]
also suffice: the last inequality gives
\(\tfrac12|F\Gamma(F,Q_h)|\le D_h+F^2P_h\).
All inequalities must hold throughout \(U\).
\end{remark}

For the natural choice \(a=\sqrt S\),
the quantities in the criterion use only \(q,S\) and their derivatives. Writing
\(K=\Delta S\) and \(T=\Gamma(S,S)\), one obtains
\begin{equation}\label{eq:harmonic-gradient-multiplier-data}
 \begin{aligned}
 h&=\tfrac12\log S,& \tau&=q/S,&
 Q_h&=\frac{T}{4S^2}+\frac{2q}{S},\\
 \Delta h&=\frac{K}{2S}-\frac{T}{2S^2},&&&
 D_h&=\frac K2-2q-\frac TS.
 \end{aligned}
\end{equation}
Under the spectral identities below, the two monotonicity signs hold,
and the remaining condition reduces to a lower bound for \(\Delta S\).

\begin{remark}[The nonharmonic correction]
\label{rem:nonharmonic-log-multiplier}
Suppose instead that \(\cL(F)=-qF\), without assuming
\(\Delta F=0\), and retain \(U=\{S>0\}\),
\eqref{eq:harmonic-log-compatibility}, and the auxiliary functions above.
The term \(2S\Delta F\) in \(\Gamma(F,S)\) changes only the defect:
\[
 \cL(aF)=a^3F\left(D_h+FQ_h\Delta F+F^2P_h
                         -\tfrac F2\Gamma(F,Q_h)\right).
\]
The gradient identity is unchanged, so the square-completion test of
\cref{rem:harmonic-alternative-tests} becomes
\[
 Q_h\ge0,\qquad D_h\ge0,\qquad P_h\ge0,\qquad
 \bigl(\Gamma(F,Q_h)-2Q_h\Delta F\bigr)^2\le16D_hP_h.
\]
It gives comparison when \(Z=\{S=0\}\) is locally
\(\cH^{N-1}\)-null, in particular for square-free homogeneous \(F\)
by \cref{lem:square-free-critical-set}. The natural choice
\(h=\frac12\log S\) must still satisfy the compatibility assumption:
now \(\Gamma(F,h)=\Delta F+qF/S\), which need not be divisible
smoothly by \(F\) across its regular zero set.
\end{remark}
